\section{Colourings}

\begin{boxdefinition}[Colouring]
    Given a graph $G = (V, E)$, a $k$-colouring of $G$ is a map $c : V \to [k]$ such that $u \sim v \implies c(u) \ne c(v)$.
\end{boxdefinition}

\subsection{$2$-Colourings of Graphs}

Let's start with a motivating question: which graphs can be 2-coloured?

Any odd cycle graph can't be 2-coloured. That is, any graph that looks like
\begin{figure}[H]
    \centering
    \begin{tikzpicture}
        \cyclegraph{3}
        \begin{scope}[xshift = 3cm]
            \cyclegraph{5}
        \end{scope}
        \begin{scope}[xshift = 6cm]
            \cyclegraph{7}
        \end{scope}
        \node at (8.25, 0.1) {$\cdots$};
    \end{tikzpicture}
\end{figure}

Let's study when odd cycles exist.

\begin{boxlemma}\label{Ch1:Lemma:Odd_Closed_Walk}
    If $G$ has an odd closed walk, then $G$ has an odd cycle.
\end{boxlemma}
Recall that a closed walk need not be a cycle: walks can contain repetitions. So even line graphs can have closed walks (start from one vertex, go to its neighbour, return).
\begin{proof}
    Suppose $G$ has an odd closed walk. Let $W$ be such a walk of minimal length. Write $W = \parenth{x_1, \ldots, x_i, \ldots, x_j, \ldots, x_{\ell}}$, with $x_{\ell} = x_1$. Suppose $W$ is not a cycle. Then there must be some $x_i, x_j$ such that $x_i = x_j$ even though $i \neq j$, ie, there must be some repetition. Then, $\parenth{x_1, \ldots, x_i, x_{j+1}, \ldots, x_{\ell}}$ and $\parenth{x_i, x_{i+1}, \ldots, x_{j}}$ are both shorter walks than $W$. Moreover, one of them must be of odd length. This contradicts minimality of the length of $W$. So $W$ must be a cycle.
\end{proof}

Here's a cool result.

\begin{boxtheorem}
    A graph $G = (V, E)$ can be $2$-coloured if and only if $G$ has no odd cycle.
\end{boxtheorem}
\begin{proof} \hfill
    \begin{description}
        \item[$\parenth{\implies}$] We've already seen that odd cycles can't be $2$-coloured. So if $G$ can be $2$-coloured, it can't have an odd cycle. We've essentially seen this from the examples above.

        \item[$\parenth{\impliedby}$] We may assume $G$ is connected, since otherwise we can colour each component separately. Fix some $v_0 \in V$. Define $c : V \to \set{1, 2}$ by
        \begin{align*}
            c(u) =
            \begin{cases}
                2 & \text{ if } \dist{v_0, u} \text{ is odd} \\
                1 & \text{ if } \dist{v_0, u} \text{ is even}
            \end{cases}
        \end{align*}
        where $\dist{\cdot, \cdot}$ denotes the shortest number of edges from one vertex to another.

        We claim that this $c$ is a valid $2$-colouring.

        Indeed, suppose it is not. Then, there must be some $u, v \in V$ such that $u \sim v$ and $c(u) = c(v)$. Now let's consider the two possible cases on the colour.

        Suppose $c(u) = c(v) = 1$ for adjacent $u, v$. Then, there is a path of even length between $v$ and $v_0$ and another path of even length between $u$ and $v_0$. Concatenating the path from $v$ to $v_0$ with the one from $v_0$ to $u$, and closing up with the edge $uv$, we get a closed walk of length even $+$ even $+ 1$. This is odd, so by \Cref{Ch1:Lemma:Odd_Closed_Walk}, $G$ has an odd cycle.

        The case $c(u) = c(v) = 2$ is identical: the two paths are now of odd length, so the closed walk they form with the edge $uv$ has length odd $+$ odd $+ 1$, which is odd again.

        Either way, $G$ has an odd cycle, contrary to hypothesis. So $c$ must be a valid $2$-colouring.
    \end{description}
\end{proof}

\subsection{Colourings of Hypergraphs}

Colouring is not really a question about graphs. All the definition above asks of $G$ is that we know which pairs of vertices may not share a colour, and nothing stops us forbidding larger sets instead.

\begin{boxdefinition}[Hypergraph]
    Given a vertex set $V$, a \textbf{hypergraph} is a collection $H$ of subsets of $V$. The sets in $H$ are called \textbf{hyperedges}.
\end{boxdefinition}

\begin{boxdefinition}[Uniformity]
    A hypergraph $H$ is \textbf{$k$-uniform} if $\abs{f} = k$ for all $f \in H$.
\end{boxdefinition}

The following is thus obvious.

\begin{boxproposition}
    A graph is a $2$-uniform hypergraph.
\end{boxproposition}

\begin{boxdefinition}
A proper colouring of a hypergraph $H$ is an assignment $c : V \to [r]$ such that $\forall f \in H$, $\abs{c(f)} \geq 2$.
\end{boxdefinition}

We say ``no edge is monochromatic''.

We can ask ourselves the following question: \textbf{``Which $3$-uniform hypergraphs are $2$-colourable?}

For each $k$, define $m(k)$ to be the smallest $m$ such that $\exists$ some $k$-uniform hypergraph with $m$ edges that is not $2$-colourable. We can show that $m(2) = 3$. But what is $m(3)$? That is, how many $3$-edges do I need to prevent $2$-colourability?

Let $H$ be a $3$-uniform hypergraph that's not $2$-colourable. Assume, further, that $H$ has at most $6$ edges (each being a subset of size $3$).

First, we'll show that we can reduce, without loss of generality, to the case where $H$ has exactly $6$ vertices. Indeed, if $H$ has vertices $x, y$ that are not in any common edge, then I can define $H'$ by ``merging'' $x$ and $y$. So $H'$ has one fewer vertex than $H$. This $H'$ cannot be $2$-colourable: if it were, then that colouring would extend to $H$ by using the colour of the merged vertex at both $x$ and $y$, which would make $H$ $2$-colourable. And when $H$ has more than $6$ vertices, such a pair always exists.

To see this, suppose $H$ has $n \geq 7$ vertices. Each edge of $H$ is a set of size $3$, and so lies over exactly $3$ pairs of vertices; since there are at most $6$ edges, at most
\begin{align*}
    6 \cdot 3 = 18
\end{align*}
of the $\binom{n}{2} \geq \binom{7}{2} = 21$ pairs of vertices are contained in a common edge. Some pair is therefore left over, and merging it brings us one vertex closer to $6$.

So now assume that $H$ has exactly $6$ edges and $6$ vertices and is not $2$-colourable. 